\section{תקציר}

\subsection{מטרת העבודה}
מטרת פרויקט זה היא לפתח מודל חיזוי דינמי מבוסס למידת מכונה \LR{(Machine Learning)} לניהול ולניטור משתני בית השורשים בחממה.

הפרויקט מתמקד ביצירת מערכת חיזוי מסוג \SoftSensor{} המסוגלת לחזות את ערכי המוליכות החשמלית (EC) והחומציות (pH) העתידיים בין מדידות ידניות דלילות, בהינתן נקודת מוצא ידועה (מדידה נוכחית).

המודל משקלל שלושה וקטורים עיקריים כדי לבצע את החיזוי:

\begin{itemize}
    \item \textbf{מצב התחלתי:} ערכי ה-EC וה-pH שנמדדו בנקודת הזמן הנוכחית.
    \item \textbf{תנאי סביבה:} נתוני אקלים חיצוני ומיקרו־אקלים בתוך החממה, לרבות קרינה, טמפרטורה, לחות, התאדות וטמפרטורת קרקע.
    \item \textbf{ממשק גידול:} נתוני השקיה, דישון, חומצה, מלחים ומועדי אירועים תפעוליים, לצד משתני מצב הצמח כגון כיסוי נוף וגיל הצמח.
\end{itemize}

\subsection{הערך היישומי}
יכולת חיזוי זו נועדה לשמש ככלי תומך החלטה עבור המגדל. היא מאפשרת לעבור מגישה תגובתית, שבה מתקנים חריגות לאחר שהתרחשו, לגישה פרואקטיבית, שבה ניתן לזהות מראש מגמות של שינויי מליחות או חומציות ולהתאים את משטר ההשקיה והדישון בהתאם.

באופן זה, המודל עשוי לסייע בשמירה על תנאים יציבים יותר בבית השורשים, בצמצום שימוש עודף במים ובדשן, ובהפחתת הסיכון לפגיעה בהתפתחות הצמח. עם זאת, בשלב הנוכחי המודל מיועד ככלי תומך החלטה ואינו מהווה מערכת בקרה אוטונומית.

\subsection{שיטת המחקר}
המחקר מתבסס על אינטגרציה של נתונים ממספר מקורות מידע, היוצרים יחד תמונת מצב רחבה של סביבת הגידול ושל פעולות הממשק בחממה:

\begin{itemize}
    \item \textbf{נתוני אקלים חיצוני:} נתוני קרינה, טמפרטורה, לחות ומשתני מזג אוויר נוספים.
    \item \textbf{נתוני חיישנים בחממה:} מדידות מיקרו-אקלים, לרבות טמפרטורה, לחות, התאדות ונתוני טמפרטורת קרקע.
    \item \textbf{נתוני תכנון גידול:} אירועי השקיה, דישון, שימוש בחומצה ומלחים, כמויות ומועדי יישום.
    \item \textbf{נתוני התפתחות הצמח:} כיסוי נוף וגיל הצמח, המשמשים לייצוג מצב הגידול וצריכת המים.
    \item \textbf{מדידות בית שורשים:} מדידות ידניות של pH ו-EC המשמשות כערכי יעד לאימון ולבדיקת המודל.
\end{itemize}

בסיס הנתונים הסופי נבנה כציר זמן אחיד ברזולוציה של 10 דקות וכלל 16{,}682 רשומות בטווח התאריכים 29.05.2025 עד 21.09.2025. עם זאת, מדידות היעד של pH ו-EC היו דלילות יחסית וכללו 109 נקודות מדידה בלבד. לכן, אחד האתגרים המרכזיים בפרויקט היה פיתוח מודל המסוגל ללמוד מתוך שילוב בין נתוני חיישנים צפופים לבין מדידות יעד ידניות ודלילות.

על בסיס נתונים אלו פותחה שרשרת מודלים דו־שלבית. בשלב הראשון נבנה מודל לחיזוי ולתיאור תנאי המיקרו-אקלים בתוך החממה. בשלב השני נבנה מודל בית שורשים מאוחד, החוזה את שינויי pH ו-EC על בסיס מצב התחלתי, תנאי סביבה, היסטוריית השקיה ודישון, ומצב הצמח.

בשלב הסופי פותח מודל היברידי המשלב בין יכולת אינטרפולציה לבין יכולת אקסטרפולציה. רכיב מבוסס עצים משמש בעיקר ללמידת קשרים מורכבים מתוך דפוסים שנצפו בנתוני העבר, ולכן מתאים לחיזוי במצבים הדומים לתנאים שנמדדו במהלך הניסוי. לצידו שולב מודל ליניארי, שמטרתו לספק תגובה יציבה יותר במצבים שבהם מופיעים אירועים תפעוליים בעלי השפעה גבוהה, כגון הוספת חומצה או שינוי חד במליחות, גם כאשר הם חורגים מהדפוסים השכיחים בנתונים. שילוב זה מאפשר למודל לחזות את שינויי ה-pH וה-EC הן על בסיס דפוסים קיימים והן על בסיס מגמות שינוי הנגזרות מהתנהגות המערכת.

\subsection{תוצאות עיקריות}
תוצאות המודל הסופי מצביעות על יכולת חיזוי טובה של משתני בית השורשים, למרות דלילות מדידות היעד. עבור pH התקבלה שגיאת MAE של כ-0.33 יחידות בבדיקת ולידציה קדימה בזמן וכ-0.29 יחידות בבדיקת הכללה. בנוסף, התקבל ערך \Rtwo{} של כ-0.93, המעיד כי המודל מצליח להסביר חלק משמעותי מהשונות בערכי החומציות. בהשוואה למודל נאיבי המבוסס על שמירת הערך האחרון שנמדד, התקבל שיפור של כ-62\% בדיוק החיזוי.

עבור EC התקבלה שגיאת MAE של כ-0.127 mS/cm בבדיקת ולידציה קדימה בזמן וכ-0.131 mS/cm בבדיקת הכללה. ערכי \Rtwo{} שהתקבלו היו כ-0.98 בבדיקת ולידציה קדימה בזמן וכ-0.95 בבדיקת הכללה. תוצאות אלו מצביעות על כך שהמודל מצליח לחזות את מגמות המליחות בצורה טובה יותר מהשוואה נאיבית, עם שיפור של כ-35\% ביחס לשימוש בערך המדידה האחרון.

עם זאת, חיזוי EC נותר מאתגר יותר מחיזוי pH, בעיקר בטווחי מוליכות נמוכים, שבהם גם סטייה אבסולוטית קטנה עשויה להתבטא כשגיאה יחסית גבוהה. לכן, אף שהתוצאות מצביעות על שיפור משמעותי ביחס לבסיס ההשוואה, נדרש המשך שיפור של רגישות המודל לשינויים עדינים בריכוזי המלחים והדשן.

באופן כללי, ממצאי הפרויקט מדגימים היתכנות לפיתוח מודל חיזוי אפקטיבי לבית השורשים בחממה. חוזקו המרכזי של המודל הוא ביכולתו לשלב נתוני אקלים, השקיה, דישון ומצב צמח לצורך חיזוי דינמי של pH ו-EC, ובכך לשמש כתשתית לכלי תומך החלטה לניהול השקיה ודישון מדויק יותר.

\newpage
